%% ====================================================================
%% CHAPTER 1: THE TRANSARC EMPIRICAL STUDY
%% ====================================================================
\chapter{The TransArc Empirical Study}
\label{ch:transarc}

\transarc recovers \sadcode trace links transitively: it first links software
architecture documentation to the architecture model (\sadsam), then links the
architecture model to source code (\samcode), and finally composes the two into
documentation-to-code links. Because the pipeline is a composition, every
\sadcode link inherits the correctness---and the errors---of both upstream
stages. This chapter dissects \transarc empirically across the five ARDoCo
benchmark projects (MediaStore, TeaStore, Teammates, BigBlueButton, and JabRef)
and establishes four findings: (i)~\sadsam is the dominant error source;
(ii)~upstream false positives are amplified by one to three orders of magnitude
through enrollment; (iii)~\sadsam links nonetheless contribute almost all
recoverable recall; and (iv)~substituting an oracle \samcode stage closes most
of the quality gap, confirming that the bottleneck lies in \sadsam rather than
\samcode.

\Cref{tab:transarc-overview} reports the end-to-end behaviour of \transarc.
File-level \fone ranges from 0.694 (MediaStore) to 0.947 (JabRef), with a
cross-project mean of 0.799. These aggregate numbers, however, conceal where the
errors originate; the remainder of this chapter decomposes them.

\begin{table}[htbp]
  \centering
  \caption{End-to-end \sadcode (\transarc) results per project: gold size, produced links, confusion counts, and \fone (with the expected/threshold values reproduced from the TLR test suite)}
  \label{tab:transarc-overview}
  \begin{tabular}{lrrrrrrrrrrr}
    \toprule
    Project & Gold & Result & TP & FP & FN & Precision & Recall & F1 & Exp P & Exp R & Exp F1 \\
    \midrule
    mediastore & 31 & 18 & 17 & 1 & 14 & 0.944 & 0.548 & 0.694 & 0.940 & 0.548 & 0.690 \\
    teastore & 27 & 20 & 20 & 0 & 7 & 1.000 & 0.741 & 0.851 & 0.999 & 0.740 & 0.850 \\
    teammates & 57 & 81 & 49 & 32 & 8 & 0.605 & 0.860 & 0.710 & 0.600 & 0.859 & 0.709 \\
    bigbluebutton & 62 & 49 & 44 & 5 & 18 & 0.898 & 0.710 & 0.793 & 0.897 & 0.709 & 0.790 \\
    jabref & 18 & 20 & 18 & 2 & 0 & 0.900 & 1.000 & 0.947 & 0.899 & 0.999 & 0.946 \\
    \bottomrule
  \end{tabular}
  \par\smallskip\footnotesize Source: \texttt{reports/TRANSARC\_EMPIRICAL\_STUDY.md}.
\end{table}


%% --------------------------------------------------------------------
\section{\sadsam as the Primary Bottleneck}
\label{sec:transarc:bottleneck}

To attribute each \transarc error to a stage, we trace every false positive
$(s, c)$---a spurious sentence-to-code link---back to the bridging model
element(s) that produced it, and classify the link by whether the fault lies in
the \sadsam stage, the \samcode stage, both, or their combination.
\Cref{tab:fp-attribution} shows the result aggregated over all projects: the
\sadsam stage alone is responsible for 93.7\% of every \transarc false positive,
versus 3.3\% for \samcode, 0.5\% for both, and 2.6\% for combination effects.

\begin{table}[htbp]
  \centering
  \caption{Attribution of every \transarc false positive to its causing stage: \sadsam dominates with 93.7\% of all false positives}
  \label{tab:fp-attribution}
  \begin{tabular}{lrr}
    \toprule
    Category & Count & Percentage \\
    \midrule
    SAD\_SAM\_CAUSED & 3440 & 93.7\% \\
    SAM\_CODE\_CAUSED & 121 & 3.3\% \\
    BOTH\_CAUSED & 17 & 0.5\% \\
    COMBINATION\_ERROR & 94 & 2.6\% \\
    Total & 3672 & 100\% \\
    \bottomrule
  \end{tabular}
  \par\smallskip\footnotesize Source: \texttt{reports/TRANSARC\_EMPIRICAL\_STUDY.md}.
\end{table}


The same asymmetry holds for false negatives. Of the recoverable misses---those
not blocked by a structural impossibility---687 of 720 (95.4\%) are caused by
missing \sadsam links, against only 33 (4.6\%) caused by missing \samcode links.
\transarc's quality is therefore almost entirely a function of \sadsam quality:
improving the architecture model or the \samcode stage in isolation cannot move
the end-to-end metric materially.

A complementary, oracle-based view confirms the ranking. If the \sadsam stage
were made perfect, the mean file-level \fone would rise from 0.794 to 0.943,
whereas a perfect \samcode stage would lift it only to 0.812; for every project
the higher-priority fix is \sadsam.

A small fraction of misses are \emph{irreducible}. \Cref{tab:theoretical-limit}
counts the gold \sadcode links for which no transitive path exists even under a
perfect upstream pipeline (i.e.\ the documentation never mentions the relevant
model element, so no composition can recover the link). Only Teammates (6.7\%)
and BigBlueButton (0.5\%) are affected; for the other three projects the
theoretical recall ceiling is 100\%. The bottleneck is thus overwhelmingly
\emph{methodological} (\sadsam errors), not a limitation of the benchmark's
ground truth.

\begin{table}[htbp]
  \centering
  \caption{Theoretical recall limit: \sadcode gold links for which no transitive path exists even with a perfect upstream pipeline. Only Teammates (6.7\%) and BigBlueButton (0.5\%) are affected}
  \label{tab:theoretical-limit}
  \begin{tabular}{lrrr}
    \toprule
    Project & Gold SAD-CODE & Theoretical Limit FNs & Unrecoverable \% \\
    \midrule
    mediastore & 59 & 0 & 0.0\% \\
    teastore & 707 & 0 & 0.0\% \\
    teammates & 8097 & 544 & 6.7\% \\
    bigbluebutton & 1529 & 8 & 0.5\% \\
    jabref & 8268 & 0 & 0.0\% \\
    \bottomrule
  \end{tabular}
  \par\smallskip\footnotesize Source: \texttt{reports/TRANSARC\_EMPIRICAL\_STUDY.md}.
\end{table}


%% --------------------------------------------------------------------
\section{Error Amplification and Cascade}
\label{sec:transarc:amplification}

A single upstream error is not a single downstream error. Enrollment expands
directory-level gold annotations into many file-level data points
(\cref{sec:distributional-inequality}), and the same expansion applies to a
system's mistakes: one wrong \sadsam link from a sentence to a component is
composed with \emph{every} code file mapped to that component, producing one
\transarc false positive per file. The result is severe error amplification.

\Cref{tab:error-amplification} quantifies this per stage and per project. A
single incorrect \sadsam link is amplified into 85.5 \transarc false positives
on average for Teammates (28~\sadsam false positives induce 2{,}395 \transarc
false positives) and into 495 for JabRef (2~\sadsam false positives induce
990). The most extreme single case is JabRef's \texttt{logic/} component: one
incorrect directory-level decision expands to 972 file-level false positives,
roughly half of that project's enrolled gold standard. By contrast, \samcode
false positives amplify far less (at most 7.5$\times$ for BigBlueButton),
because a component typically maps to many files but a sentence maps to few
components.

\begin{table}[htbp]
  \centering
  \caption{Error amplification: each upstream false positive cascades into many \transarc false positives. The mean amplification factor reaches 85.5 for Teammates and 495 for JabRef}
  \label{tab:error-amplification}
  \begin{tabular}{lrrrrrr}
    \toprule
    Project & \sadsam FPs & Induced \transarc FPs & Amp. (\sadsam) & \samcode FPs & Induced \transarc FPs & Amp. (\samcode) \\
    \midrule
    mediastore & 1 & 1 & 1.0 & 0 & 0 & 0.0 \\
    teastore & 0 & 0 & 0.0 & 0 & 0 & 0.0 \\
    teammates & 28 & 2395 & 85.5 & 0 & 0 & 0.0 \\
    bigbluebutton & 5 & 71 & 14.2 & 18 & 135 & 7.5 \\
    jabref & 2 & 990 & 495.0 & 1 & 4 & 4.0 \\
    \bottomrule
  \end{tabular}
  \par\smallskip\footnotesize Source: \texttt{reports/TRANSARC\_EMPIRICAL\_STUDY.md}.
\end{table}


This cascade is the mechanism behind the ranking instability discussed in
\cref{ch:bias}: because file-level scoring weights each annotator decision by
its directory's file count, a handful of upstream mistakes on large components
can dominate the metric, and the amplification factor differs by two orders of
magnitude across projects.

%% --------------------------------------------------------------------
\section{Actual Contribution of \sadsam}
\label{sec:transarc:contribution}

Having shown that \sadsam dominates the errors, we now quantify what \sadsam
\emph{contributes}. \Cref{tab:sadsam-contribution} lists, per project, the
\sadsam true and false positives that seed the transitive pipeline together
with the \sadcode links they ultimately yield. The relationship is highly
non-linear: TeaStore's 20~correct \sadsam links expand into 501~\sadcode true
positives, and JabRef's 18~correct links expand into 8{,}268. Correct \sadsam
links are the engine of recall; the entire \transarc result is built on a few
dozen architecture-level decisions per project.

\begin{table}[htbp]
  \centering
  \caption{Actual contribution of \sadsam to \transarc: the \sadsam true/false positives that seed the transitive pipeline and the resulting \sadcode links, per project}
  \label{tab:sadsam-contribution}
  \begin{tabular}{lrrrrr}
    \toprule
    Project & SAD-SAM TPs & SAD-SAM FPs & SAD-CODE TPs & SAD-CODE FPs & SAD-CODE Total \\
    \midrule
    mediastore & 17 & 1 & 25 & 1 & 26 \\
    teastore & 20 & 0 & 501 & 0 & 501 \\
    teammates & 49 & 32 & 7307 & 2395 & 9702 \\
    bigbluebutton & 44 & 5 & 1287 & 282 & 1569 \\
    jabref & 18 & 2 & 8268 & 994 & 9262 \\
    \bottomrule
  \end{tabular}
  \par\smallskip\footnotesize Source: \texttt{reports/SAD\_SAM\_ACTUAL\_CONTRIBUTION.md}.
\end{table}


The flip side is the recall that \emph{better} \sadsam would unlock.
\Cref{tab:sadsam-tp-gain} estimates, for each \sadsam false negative, the
\sadcode true positives that would become recoverable if that link were fixed.
The average gain per recovered \sadsam false negative reaches 30.8 \sadcode true
positives for Teammates and 29.4 for TeaStore, with single-link gains as high as
71. \Cref{tab:sadsam-coverage} turns this into a recall ceiling: recovering all
\sadsam false negatives would address 100\% of TeaStore's \sadcode false
negatives, 79.4\% of MediaStore's, and 74.0\% of BigBlueButton's, but only
31.1\% of Teammates' (the remainder being the theoretical-limit misses of
\cref{tab:theoretical-limit}). JabRef has no \sadsam false negatives and thus no
\sadsam-recoverable recall.

\begin{table}[htbp]
  \centering
  \caption{\sadsam true-positive gain: \sadcode true positives that would become recoverable if each \sadsam false negative were fixed}
  \label{tab:sadsam-tp-gain}
  \begin{tabular}{lrrrr}
    \toprule
    Project & SAD-SAM FNs & Total Potential New SAD-CODE TPs & Max Single-Link Gain & Avg Gain/FN \\
    \midrule
    mediastore & 14 & 27 & 3 & 1.9 \\
    teastore & 7 & 206 & 64 & 29.4 \\
    teammates & 8 & 246 & 71 & 30.8 \\
    bigbluebutton & 18 & 195 & 22 & 10.8 \\
    jabref & 0 & 0 & 0 & 0.0 \\
    \bottomrule
  \end{tabular}
  \par\smallskip\footnotesize Source: \texttt{reports/SAD\_SAM\_TP\_GAIN\_STUDY.md}.
\end{table}


\begin{table}[htbp]
  \centering
  \caption{Recoverability of \sadcode false negatives through \sadsam false-negative recovery (upper bound on recall gain)}
  \label{tab:sadsam-coverage}
  \begin{tabular}{lrrr}
    \toprule
    Project & SAD-CODE FNs & Recoverable via SAD-SAM FN Recovery & Coverage \\
    \midrule
    mediastore & 34 & 27 & 79.4\% \\
    teastore & 206 & 206 & 100.0\% \\
    teammates & 790 & 246 & 31.1\% \\
    bigbluebutton & 242 & 179 & 74.0\% \\
    jabref & 0 & 0 & 0.0\% \\
    \bottomrule
  \end{tabular}
  \par\smallskip\footnotesize Source: \texttt{reports/SAD\_SAM\_TP\_GAIN\_STUDY.md}.
\end{table}


Together these tables establish the central asymmetry of \transarc: \sadsam is
simultaneously the source of almost all errors \emph{and} the lever for almost
all recoverable recall. Effort spent on the \samcode stage yields little; effort
spent on \sadsam yields nearly everything.

%% --------------------------------------------------------------------
\section{The \samcode Cascade}
\label{sec:transarc:samcode}

\samcode is the second link of the cascade. Although it causes only 3.3\% of
\transarc false positives overall, its errors are concentrated in
BigBlueButton, where the architecture model maps several components to large,
loosely related code regions. \Cref{tab:samcode-cascade} shows that
BigBlueButton's \samcode false positives induce 135 \sadcode false
positives---47.9\% of that project's \transarc false positives---whereas for
MediaStore, TeaStore, and Teammates the \samcode stage induces none.

\begin{table}[htbp]
  \centering
  \caption{\samcode stage as the second link of the \transarc cascade: \samcode false positives and the \sadcode false positives they induce, per project}
  \label{tab:samcode-cascade}
  \begin{tabular}{lrrrrr}
    \toprule
    Project & \samcode FPs & Induced \sadcode FPs & Amp. (sents/FP) & \sadcode FPs from \samcode & \% of all \transarc FPs \\
    \midrule
    mediastore & 0 & 0 & 0.0 & 0 & 0.0\% \\
    teastore & 0 & 0 & 0.0 & 0 & 0.0\% \\
    teammates & 0 & 0 & 0.0 & 0 & 0.0\% \\
    bigbluebutton & 23 & 135 & 5.9 & 135 & 47.9\% \\
    jabref & 1 & 4 & 4.0 & 4 & 0.4\% \\
    \bottomrule
  \end{tabular}
  \par\smallskip\footnotesize Source: \texttt{reports/SAM\_CODE\_CASCADE.md}.
\end{table}


The lesson mirrors the \sadsam analysis: \samcode errors matter only where the
model-to-code mapping is many-to-many and the offending component carries many
files. Where the mapping is clean, \samcode is effectively transparent and the
end-to-end quality is decided entirely upstream.

%% --------------------------------------------------------------------
\section{S12C/S12E versus \transarc}
\label{sec:transarc:s12c}

To isolate the contribution of each stage we compare three configurations:
end-to-end \transarc (TA), and two variants that replace the \samcode stage with
an oracle---S12C (oracle \samcode component mapping) and S12E (oracle \samcode
extended). If the bottleneck were \samcode, oracle substitution would barely
change the result; if it were \sadsam, the result would improve only to the
extent that \sadsam is correct.

\Cref{tab:s12c-four-level} first shows that even within \transarc the choice of
aggregation level reorders the projects: JabRef is best on file-level \fone
(0.943) but worst on decision-level \fone (0.394), because almost all of its
gold is a single huge directory block (\cref{ch:bias}). This motivates reporting
multiple levels rather than file-level \fone alone.

\begin{table}[htbp]
  \centering
  \caption{Four-level \sadcode evaluation of \transarc: file-, decision-, component-, and inverse-expansion-weighted \fone}
  \label{tab:s12c-four-level}
  \begin{tabular}{lrrrr}
    \toprule
    Project & File \fone & Decision \fone & Component \fone & Weighted \fone \\
    \midrule
    mediastore & 0.588 & 0.568 & 0.642 & 0.568 \\
    teastore & 0.829 & 0.824 & 0.828 & 0.824 \\
    teammates & 0.821 & 0.564 & 0.648 & 0.550 \\
    bigbluebutton & 0.831 & 0.629 & 0.574 & 0.623 \\
    jabref & 0.943 & 0.394 & 0.880 & 0.394 \\
    Average & 0.802 & 0.596 & 0.714 & 0.592 \\
    \bottomrule
  \end{tabular}
  \par\smallskip\footnotesize Source: \texttt{reports/S12C\_VS\_TRANSARC.csv} (\transarc rows). Identical \sadcode trace links scored at four aggregation levels. JabRef is best on File \fone (0.943) yet worst on Decision \fone (0.394), illustrating the ranking instability of file-level scoring.
\end{table}


\Cref{tab:transarc-dashboard} places \transarc beside the S12C oracle-\samcode variant.
Substituting an oracle \samcode raises the mean file-level \fone from 0.802 to
0.928 and the mean decision-level \fone from 0.596 to 0.817. The gap that
remains under S12C is precisely the \sadsam error that no \samcode improvement
can repair---confirming, from a third independent angle, that \sadsam is the
binding constraint on \transarc. The MediaStore jump is the most dramatic
(decision \fone 0.568~$\to$~0.926), while Teammates and BigBlueButton improve
least, consistent with their larger share of \sadsam-caused and combination
errors.

\begin{table}[htbp]
  \centering
  \caption{Evaluation dashboard: \transarc versus the \samcode-oracle (S12C) variant, scored at multiple aggregation levels with cross-project averages}
  \label{tab:transarc-dashboard}
  \begin{tabular}{lrrrrr}
    \toprule
    Project & TA File \fone & TA Dec. \fone & S12C File \fone & S12C Dec. \fone & S12C Comp. \fone \\
    \midrule
    mediastore & 0.588 & 0.568 & 0.929 & 0.926 & 0.986 \\
    teastore & 0.829 & 0.824 & 0.977 & 0.971 & 0.974 \\
    teammates & 0.821 & 0.564 & 0.870 & 0.589 & 0.549 \\
    bigbluebutton & 0.831 & 0.629 & 0.867 & 0.661 & 0.634 \\
    jabref & 0.943 & 0.394 & 0.999 & 0.938 & 0.898 \\
    Average & 0.802 & 0.596 & 0.928 & 0.817 & 0.808 \\
    \bottomrule
  \end{tabular}
  \par\smallskip\footnotesize Source: \texttt{reports/S12C\_VS\_TRANSARC.csv}. ``TA'' = end-to-end \transarc; ``S12C'' substitutes the oracle \samcode link. The spread between File and Decision \fone quantifies enrollment-inflation bias; the TA$\to$S12C jump isolates the \sadsam bottleneck.
\end{table}


%% --------------------------------------------------------------------
\section{Summary}
\label{sec:transarc:summary}

The empirical study yields a coherent picture. \transarc's end-to-end quality is
governed by the \sadsam stage: \sadsam causes 93.7\% of false positives and
95.4\% of recoverable false negatives, carries nearly all recoverable recall,
and---once replaced downstream by an oracle \samcode---still bounds the
achievable \fone. Enrollment amplifies each upstream error by up to three orders
of magnitude, so a tiny number of architecture-level mistakes dominates the
file-level metric and destabilises cross-project rankings. These two
observations---an upstream bottleneck and an enrollment amplifier---motivate the
benchmark-bias and evaluation-metric analysis of \cref{ch:bias}.
